\chapter{Conclusao}
\label{chap_conclusao}

O trabalho atingiu o objetivo de implementar e analisar automatos celulares elementares de Wolfram com persistencia completa dos resultados (dados de entrada, dados de saida e figuras). A organizacao dos dados em formato padronizado, incluindo CSV, permitiu rastrear de forma direta cada experimento e seus parametros.

Do ponto de vista qualitativo, as oito regras selecionadas apresentaram comportamento coerente com as quatro classes de Wolfram. Do ponto de vista quantitativo, as metricas extraidas dos dados em CSV confirmaram a separacao entre classes de baixa atividade (homogenea e periodica) e classes com atividade mais intensa (caotica e complexa).

Na comparacao de escala, verificou-se que algumas regras sao mais sensiveis ao aumento do tamanho da malha do que outras. Em especial, a Regra 45 manteve densidade media praticamente constante, enquanto as regras 30, 54 e 110 apresentaram reducao de densidade media para \(tamanho=200\).

Como continuidade, recomenda-se:
\begin{itemize}
    \item incluir metricas adicionais de complexidade (entropia e compressibilidade);
    \item automatizar a classificacao das regras com base em descritores numericos;
    \item expandir a analise para condicoes iniciais alternativas.
\end{itemize}
